\chapter{\selectlanguage{greek}Θεωρητικό υπόβαθρο}

Για την ανάπτυξη του μοντέλου αξιοποιήθηκαν έννοιες και
τεχνολογίες από διαφορετικούς τομείς. Στο παρόν κεφάλαιο
παρουσιάζεται το θεωρητικό πλαίσιο που υποστηρίζει αυτή την
εργασία.

\section{Υπερτουρισμός}

Υπερτουρισμός \en{(overtourism)} είναι η κατάσταση κατά την οποία
ο αριθμός επισκεπτών σε έναν προορισμό υπερβαίνει τη $``$φέρουσα
ικανότητα$"$ \en{(carrying capacity)} του, επηρεάζοντας αρνητικά
τους κατοίκους, το περιβάλλον και τη ποιότητα εμπειρίας των
επισκεπτών. Ενώ ο απλός τουρισμός θεωρείται θετικός οικονομικός
πυλώνας που ενισχύει την πολιτισμική ανταλλαγή, ο υπερτουρισμός
απειλεί την αυθεντικότητα και τη ποιότητα ζωής των κατοίκων.

Όταν μια περιοχή δέχεται περισσότερους τουρίστες από όσους μπορεί
να καλύψει, τότε οι υποδομές της καταρρέουν, είτε αυτές έχουν να
κάνουν με νερό, με αποχετεύση, με ηλεκτρικό ρεύμα, με διαδίκτυο,
ή με συγκοινωνίες. Οι μόνιμοι κάτοικοι συνήθως εκτοπίζονται από
κεντρικές περιοχές λόγω αύξησης των ενοικίων. Οι τουρίστες από
την άλλη, δεν απολαμβάνουν τον προορισμό, λόγω συνωστισμού, ουρών
και απώλεια της αυθεντικότητας του. Το αποτύπωμα του υπερτουρισμού
γίνεται ιδιαίτερα αντιληπτό και στο περιβάλλον, καθώς η 
ανεξέλεγκτη ρύπανση παντός τύπου καταστρέφει οικοσυστήματα και
αλλοιώνει φυσικά τοπία. Σε βάθος χρόνου, υπερτουρισμός συχνά
σημαίνει ότι μια περιοχή εξαρτάται αποκλειστικά από τον τουρισμό,
πλήττοντας άλλες οικονομικές δραστηριότητες. Στην Ελλάδα, το
φαινόμενο εντοπίζεται έντονα σε προορισμούς όπως η Σαντορίνη και
η Μύκονος, όπου η ισορροπία μεταξύ κέρδους και βιωσιμότητας έχει
διαταραχθεί. Την πρωτιά ωστόσο την κατέχει η Ζάκυνθος, ως το νησί
με τη μεγαλύτερη αναλογία τουριστών ανά μόνιμο κάτοικο στην Ευρώπη.
\\(\en{https://www.kathimerini.gr/k/travel/563622064/posoys-toyristes-mporei-na-antexei-i-zakynthos/})

\section{\en{Agent Based Models (ABMs)}}

\subsection{Τι είναι τα \en{ABMs}}

Ένα \en{Agent Based Model (ABM)} είναι μια υπολογιστική μέθοδος
προσομοίωσης συστημάτων μέσω αυτόνομων οντοτήτων, γνωστές και ως
\en{agents} (πράκτορες). Οι πράκτορες έχουν κανόνες συμπεριφοράς
και αλληπιδρούν τόσο μεταξύ τους όσο και με το περιβάλλον. 

Σε ένα μοντέλο μπορούμε να έχουμε παραπάνω από έναν τύπους
πρακτόρων. Για παράδειγμα, σε ένα μοντέλο μπορεί έχουμε κλάση
$``$κάτοικος$"$ και κλάση $``$τουρίστας$"$.
Κάθε κλάση πρακτόρων ορίζει τα δικά της χαρακτηριστικά και κανόνες
συμπεριφοράς για τους πρακτόρες που ανήκουν σε αυτή. Έτσι, οι
\en{agents} τύπου $``$κάτοικος$"$ έχουν διαφορετικά χαρακτηριστικά
και κανόνες από τους \en{agents} που ανήκουν στους τουρίστες.
Κάθε πράκτορας μπορεί να αλληλεπιδρά τόσο με το περιβάλλον, όσο
και με άλλους πράκτορες που ανήκουν στην ίδια ή διαφορετική κλάση
με αυτόν, ανάλογα με τους προβλεπόμενους κανόνες της κλάσης του.
Δηλαδή ένας \en{agent} τύπου $``$τουρίστας$"$ μπορεί να αλληλεπιδρά
τόσο με το περιβάλλον όσο και με κατοίκους ή με άλλους τουρίστες.

Ένα \en{Agent Based Model} μάς επιτρέπει να εκτελέσουμε
προσωμοιώσεις φαινομένων οι οποίες εξελίσσονται με το χρόνο.
Για την ακρίβεια, η προσομοίωση $``$τρέχει$"$ σε διακριτά βήματα
(\en{ticks} ρολογιού), όπου μετά από κάθε βήμα το μοντέλο
μεταβαίνει σε μια καινούρια επόμενη κατάσταση.

Η μαγεία των \en{ABMs} αναδεικνύεται μέσα από την αναδυόμενη
συμπεριφορά που συχνά παρουσιάζουν. Με άλλα λόγια, το σύστημα ως
σύνολο συχνά εκδηλώνει ιδιότητες που δεν προβλέπονται από τους
κανόνες του κάθε πράκτορα ξεχωριστά.

\subsection{Το \en{Mesa} ως \en{Framework} Μοντελοποίησης}

Το \en{Mesa} είναι ένα \en{framework} ανοιχτού κώδικα της
\en{Python} για Μοντελοποίηση Βάσει Πρακτόρων. Αν και συχνά
αναφέρεται ως βιβλιοθήκη, στην πραγματικότητα αποτελεί ένα
ολοκληρωμένο \en{framework}, καθώς επιβάλει μια συγκεκριμένη
αρχιτεκτονική δόμησης των μοντέλων. Σχεδιάστηκε ως μια σύγχρονη
εναλλακτική σε παλαιότερα εργαλεία (όπως το \en{NetLogo} ή το
\en{RePast}), δίνοντας έμφαση στην ευελιξία και την αξιοποίηση
του οικοσυστήματος της \en{Python} για επιστημονική ανάλυση
δεδομένων.

Η ισχύς του \en{Mesa} έγκειται στην αρθρωτή του δομή
\en{(modular design)}, επειδή παρέχει ανεξάρτητα υποστυστήματα
τα οποία οι ερευνητές μπορούν να συνδυάσουν ανάλογα με τις ανάγκες
τους. Αυτή η αρχιτεκτονική επιτρέπει τον σαφή διαχωρισμό μεταξύ
της λογικής του μοντέλου, της διαχείρισης του χρόνου και της
οπτικοποίησης των αποτελεσμάτων.

Τα βασικά χαρακτηριστικά του \en{framework} είναι τέσσερα. Πρώτον,
η διαχείριση χρόνου και χώρου, μέσω έτοιμων μηχανισμών για τον
προγραμματισμό ενεργειών \en{(Schedulers)}, όπως η ταυτόχρονη ή η
τυχαία ενεργοποίηση πρακτόρων, καθώς και ποικίλων τύπων πλεγμάτων
\en{(Grids)} για την κίνηση των οντοτήτων. Δεύτερον, η οπτικοποίηση
σε πραγματικό χρόνο, μέσω εργαλείων για τη δημιουργία διαδραστικών
διεπαφών μέσω \en{browser} ή \en{Jupyter Notebooks}, που επιτρέπουν
στον χρήστη να παρατηρεί τη δυναμική της προσομοίωσης και να
πειραματίζεται με παραμέτρους εύκολα και γρήγορα. Τρίτον, η
ενσωμάτωση στο \en{Python ecosystem}, η οποία συνδέει το \en{Mesa}
άψογα με βιβλιοθήκες όπως η \en{Pandas} για την ανάλυση δεδομένων
και η \en{Matplotlib} για τη δημιουργία στατικών γραφημάτων, ενώ η
συλλογή δεδομένων \en{(Data Collection)} γίνεται αυτόματα,
διευκολύνοντας τη στατιστική επεξεργασία μετά την προσομοίωση.
Τέλος, η επεκτασιμότητα, καθώς το \en{framework} υποστηρίζει
εξειδικευμένες επεκτάσεις, με σημαντικότερη το \en{Mesa-Geo}, το
οποίο επιτρέπει την ενσωμάτωση γεωγραφικών δεδομένων \en{(GIS)} για
τη μοντελοποίηση πρακτόρων σε πραγματικούς χάρτες.

\section{Γεωχωρικά Δεδομένα \en{OpenStreetMap}}

\subsection{\en{OpenStreetMap (OSM)}}

Το \en{OpenStreetMap (OSM)} αποτελεί μια παγκόσμια, συνεργατική
βάση γεωγραφικών δεδομένων που λειτουργεί υπό τις αρχές του
ανοιχτού περιεχομένου. Συχνά αναφέρεται ως η «\en{Wikipedia} των
χαρτών», καθώς συντηρείται από μια κοινότητα εθελοντών που
συνεισφέρουν δεδομένα για οδικά δίκτυα, κτίρια, φυσικά
χαρακτηριστικά και σημεία ενδιαφέροντος \en{(Points of Interest - 
POI)}. Η ανοιχτή φύση των δεδομένων του, σε συνδυασμό με την
υψηλή λεπτομέρεια που προσφέρει, το καθιστά το βασικό εργαλείο για
την ανάπτυξη εφαρμογών \en{GIS} και προσομοιώσεων που απαιτούν
ρεαλιστική χωρική πληροφορία.

\subsection{Η Μορφή Δεδομένων \en{Shapefile}}

Τα γεωχωρικά δεδομένων του \en{OSM} διατίθενται συχνά σε μορφή
\en{Shapefile (.shp)}, η οποία αποτελεί ένα βιομηχανικό πρότυπο
για την αποθήκευση διανυσματικών δεδομένων \en{(vector data)}.
Ένα \en{Shapefile} δεν είναι ένα μεμονομένο αρχείο, αλλά ένα
σύνολο αρχείων με κοινό όνομα και διαφορετικές επεκτάσεις.
Ενδεικτικά:

\begin{itemize}
\item \en{.shp}: Περιέχει τη γεωμετρία των οντοτήτων (σημεία, 
γραμμές, πολύγωνα).
\item \en{.dbf}: Αποθηκεύει τα περιγραφικά χαρακτηριστικά
\en{(attributes)} κάθε γεωμετρίας σε μορφή πίνακα.
\item \en{.shx}: Λειτουργεί ως ευρετήριο \en{(index)} για τη
σύδνεση γεωμετρίας και δεδομένων
\end{itemize}

\subsection{Συστήματα Συντεταγμένων και Προβολές}

Μια κρίσιμη πτυχή της διαχείρισης γεωχωρικών δεδομένων είναι η
κατανόηση των Συστημάτων Αναφοράς Συντεταγμένων \en{(Coordinate
Reference Systems - CRS)}.

Το διεθνές γεωδαιτικό σύστημα που χρησιμοποιείται από τα \en{GPS}
ονομάζεται \en{WGS84 (EPSG:4326)}. Αντιμετωπίζει τη Γη ως
σφαιροειδές και ορίζει τις θέσεις σε μοίρες (γεωγραφικό μήκος και
πλάτος). Ενώ είναι ιδανικό για παγκόσμιο εντοπισμό, παρουσιάζει
δυσκολίες στον υπολογισμό ευκλείδιων αποστάσεων (π.χ. μέτρηση
σε μέτρα) λόγω της καμπυλότητας της Γης.

Το ΕΓΣΑ87 \en{(EPSG:2100)} ή αλλιώς Ελληνικό Γεωδαιτικό Σύστημα
Αναφοράς 1987 είναι ένα προβολικό σύστημα σχεδιασμένο ειδικά για
τον ελλαδικό χώρο. $``$Μεταφράζει$"$ την καμπύλη επιφάνεια σε ένα
επίπεδο καρτεσιανό σύστημα \en{(X, Y)} σε μέτρα. Η χρήση
προβολικών συστημάτων είναι απαραίτητη σε προσομοιώσεις πρακτόρων,
καθώς επιτρέπει τον ακριβή υπολογισμό αποστάσεων και ταχυτήτων με
τη χρήση απλής γεωμετρίας.

\subsection{Εργαλεία Διαχείρισης \en{(GeoPandas)}}

Για την προγραμματαιστική επεξεργασία των παραπάνω δεδομένων, η
βιβλιοθήκη \en{GeoPandas} αποτελεί το βασικό εργαλείο στο
οικοσύστημα της \en{Python}. Επεκτείνει τις δυνατότητες της
βιβλιοθήκης \en{Pandas} προσθέτοντας υποστήριξη για γεωχωρικές
λειτουργίες μέσω της δομής \en{GeoDataFrame}. Επιτρέπει το
χωρικό φιλτράρισμα δεδομένων \en{(Spacial Join/Clipping)} για τον
περιορισμό της ανάλυσης σε συγκεκριμένα γεωγραφικά όρια
\en{(Bounding Box)}, καθώς και τον εύκολο μετασχηματισμό μεταξύ
διαφορετικών συστημάτων συντεταγμένων \en{(Reprojection)}. Τέλος,
επιτρέπει τη διαχείριση θεματικών επιπέδων \en{(Layers)}, όπου
κάθε επίπεδο αντιπροσωπεύει μια διαφορετική κατηγορία πληροφορίας,
όπως το οδικό δίκτυο, τα σημεία ενδιαφέροντος ή οι χρήσεις γης.


\section{Θεωρία Γράφων για Οδικά Δίκτυα}

\subsection{Βασικές Έννοιες}

Στο πλαίσιο της μαθηματικής μοντελοποίησης, ένας γράφος
\en{(graph) G=(V,E)} ορίζεται ως μια δομή που αποτελείται από ένα
σύνολο κορυφών ή κόμβων \en{V (vertices)} και ένα σύνολο ακμών
\en{E (edges)}, οι οποίες συνδέουν ζεύγη κορυφών. Αν οι ακμές του
γράφου έχουν συγκεκριμένη φορά (π.χ. μονόδρομοι) τότε το γράφημα
λέγεται Κατευθυνόμενο. Αντιθέτως, αν η σύνδεση μεταξύ των κόμβων
είναι αμφίδρομη τότε το γράφημα είναι Μη Κατευθυνόμενο. Υπάρχουν
επίσης Γράφοι με Βάρη \en{(Weighted Graphs)}, όπου σε κάθε ακμή
του γράφου ανατίθεται μια αριθμητική τιμή (βάρος), η οποία μπορεί
να αντιπροσωπεύει το φυσικό μήκος, τον χρόνο διέλευσης ή το
κόστος μετακίνησης.

\subsection{Μοντελοποίηση Οδικών Δικτύων}

Η αναπαράσταση ενός οδικού δικτύου μέσω της θεωρίας γράφων
επιτρέπει την εφαρμογή αλγορίθμων βελτιστοποίησης για τη
μετακίνηση πρακτόρων στον χώρο. Στη μοντελοποίηση αυτή, οι κόμβοι
αντιπροσωπεύουν κρίσιμα σημεία του δικτύου, όπως διασταυρώσεις
δρόμων ή τέρματα αδιεξόδων, ενώ οι ακμές αντιπροσωπεύουν τα οδικά
τμήματα που συνδέουν τους κόμβους. Επιπλέον, τα βάρη των ακμών
καθορίζουν τη λήψη αποφάσεων των πρακτόρων. Συνήθως
χρησιμοποιείται ευκλείδια απόσταση, αλλά μπορούν να ενσωματωθούν
και ποιοτικά χαρακτηριστικά, όπως η ελκυστικότητα μιας διαδρομής
για έναν πεζό.

Η χρήση γράφων μετατρέπει το πρόβλημα της πλοήγησης σε πρόβλημα
εύρεσης συντομότερης διαδρομής \en{(pathfinding)}, το οποίο
επιλύεται με κλασικούς αλγορίθμους όπως ο \en{Dijkstra} ή ο
\en{A}*.

\subsection{Η βιβλιοθήκη \en{NetworkX}}

Για την ανάλυση και τη διαχείριση αυτών των δομών στην \en{Python},
το \en{NetworkX} αποτελεί την τυπική βιβλιοθήκη για αυτό το σκοπό.
Παρέχει ένα ευρύ φάσμα για τη δημιουργία και τροποποίηση σύνθετων
δικτύων (γράφων) και τον έλεγχο της συνδεσιμότητας
\en{(connectivity)}, διασφαλίζοντας ότι υπάρχει μονοπάτι μεταξύ
των κόμβων του δικτύου. Επιπροσθέτως, επιτρέπει την εξαγωγή
τοπολογικών χαρακτηριστικών (π.χ. βαθμός κόμβου, κεντρικότητα), τα
οποία βοηθούν στην κατανόηση της δομής του αστικού περιβάλλοντος.
